\section{2017秋}


\subsection{微分積分}
\prob{
  次の極限値を求めなさい。
  \begin{enumerate}[label=(\arabic*), ref=(\arabic*), itemsep=5pt]
    \item \dm{\lim_{x\to\infty} \bigl(\sqrt{x^2 + 2x - 3} - x - 3\bigr)}
    \item \dm{\lim_{x\to 0} (\sin x + \cos x)^{1/x}}
  \end{enumerate}
}


\prob{
  次の重積分について、以下の問に答えなさい。
  $e$は自然対数の底（ネイピア数）である。
  \begin{gather}
    \iint_{D} y^2e^{xy}\,dxdy,\quad D = \Dset{(x,y) \relmiddle 0<x\leq y\leq 2x,\,1\leq xy\leq 4}
  \end{gather}
  \begin{enumerate}[label=(\arabic*), ref=(\arabic*), itemsep=5pt]
    \item 積分領域$D$を図示しなさい。
    \item \label{item:2017a-cal-2-2}次の変数変換のヤコビアン\dm{\frac{\partial(s,t)}{\partial(x,y)}}を求めなさい。
    \begin{align*}
        \left\{\begin{aligned}
          s = xy \\
          t = \frac{y}{x}
        \end{aligned}\right.
        \quad (x>0,\,y>0)
    \end{align*}
    \item 問\ref*{item:2017a-cal-2-2}の変数変換によって積分領域$D$が
    $(s,t)$平面の領域$\bar{D}$へうつされるとする。
    上記の重積分を領域$\bar{D}$での重積分に書き換えなさい。
    \item この重積分を計算しなさい。
  \end{enumerate}
}
% (3)で使うヤコビアンは$\frac{\partial (x,y)}{\partial(s,t)}$？

\prob{
  次の微分方程式を解きなさい。
  \begin{enumerate}[label=(\arabic*), ref=(\arabic*), itemsep=5pt]
    \item \dm{1 - y^2 - 2xyy' = 0}
    \item \dm{x^2y' + y - 2xy - x^2 = 0}
  \end{enumerate}
}

\newpage
\subsection{線形代数}

\prob{
  $\gro$を$\gro^3 = 1$を満たすある数とするとき、下記の行列式を求めよ。
  \begin{enumerate}[label=(\arabic*), ref=(\arabic*), itemsep=10pt]
    \item \dm{
      a = \vmat{
        1 & \gro & \gro^2 \\
        \gro & \gro^2 & \gro^3 \\
        \gro^2 & \gro^3 & 1
      }
    }
    \item \dm{
      b = \vmat{
        1 & \gro & \gro^2 & \gro^3 \\
        \gro & \gro^2 & \gro^3 & 1 \\
        \gro^2 & \gro^3 & 1 & \gro \\
        \gro^3 & 1 & \gro & \gro^2
      }
    }
  \end{enumerate}
}

\prob{
  \begin{enumerate}[label=(\arabic*), ref=(\arabic*), itemsep=0pt]
    \item 下記の行列の階数を求めよ。
    \begin{gather}
      \bA = \pmat{
        1 & 2 & -1 \\
        2 & -1 & 2 \\
        3 & 1 & 1 \\
        4 & 3 & 0 \\
        1 & -3 & 3
      }
    \end{gather}
    \item ある定数$a$を含む下記の連立方程式を解け。
    \begin{gather}
      \pmat{
        1 & 2 & -1 \\
        2 & -1 & 2 \\
        4 & 3 & 0
      }\pmat{
        x \\ y \\ z
      } = \pmat{
        4 \\ 2 \\ a
      }
    \end{gather}
  \end{enumerate}
}

\prob{
  \begin{gather}
    \bA = \pmat{
      \disp\frac{5}{2} & 1 \\
      \\
      1 & 1
    }
  \end{gather}
  \begin{enumerate}[label=(\arabic*), ref=(\arabic*), itemsep=0pt]
    \item この行列$\bA$の固有値と固有ベクトルを求めよ。
    \item $\bA^n$を求めよ。ただし、$n$は任意の自然数とする。
  \end{enumerate}
}
